\subsubsection{Ângulos de Euler}
\label{sec:angulos_euler}

Assim como a matriz de cossenos diretores, a orientação de um referencial $B$
em relação a outro $A$ pode ser descrita por ângulos de Euler.
Conforme \cite{wie2008space}, considera-se a sequência de rotações
$C_1(\phi_1) \leftarrow C_2(\phi_2) \leftarrow C_3(\phi_3)$ que leva o
referencial inercial $A$ até o sistema do corpo $B$. Introduzem-se dois
referenciais intermediários $A'$ e $A''$:

\begin{equation}
C_3(\phi_3): A' \leftarrow A,\qquad
C_2(\phi_2): A'' \leftarrow A',\qquad
C_1(\phi_1): B \leftarrow A'',
\end{equation}

de modo que a orientação final é dada por

\begin{equation}
\mathbf{C}^{B/A}
=
C_1(\phi_1)\,C_2(\phi_2)\,C_3(\phi_3).
\end{equation}

As velocidades angulares associadas a cada rotação elementar são
\begin{equation}
\vec{\omega}_{A'/A} = \dot{\theta}_3\,\vec{a}_3
,\qquad
\vec{\omega}_{A''/A'} = \dot{\theta}_2\,\vec{a}'_2
,\qquad
\vec{\omega}_{B/A''} = \dot{\theta}_1\,\vec{b}_1,
\end{equation}
onde $\{\vec{a}_1,\vec{a}_2,\vec{a}_3\}$ é a base de $A$ e
$\{\vec{b}_1,\vec{b}_2,\vec{b}_3\}$ a base de $B$.
O vetor velocidade angular total de $B$ em relação a $A$ é a soma das três
contribuições~\cite{wie2008space}:
\begin{equation}
\vec{\omega}^{B/A}
=
\vec{\omega}_{B/A''}
+
\vec{\omega}_{A''/A'}
+
\vec{\omega}_{A'/A}
=
\dot{\theta}_1\,\vec{b}_1
+
\dot{\theta}_2\,\vec{a}''_2
+
\dot{\theta}_3\,\vec{a}'_3.
\label{eq:omega_soma_contrib}
\end{equation}

Escrevendo $\vec{\omega}^{B/A}$ nas componentes do corpo,
$\vec{\omega}^{B/A} = \omega_1 \vec{b}_1 + \omega_2 \vec{b}_2 + \omega_3 \vec{b}_3$,
e projetando o lado direito de \eqref{eq:omega_soma_contrib} nos eixos de $B$,
Wie obtém a relação matricial~\cite{wie2008space}:
\begin{equation}
\begin{bmatrix}
\omega_1 \\[2pt] \omega_2 \\[2pt] \omega_3
\end{bmatrix}
=
\begin{bmatrix}
1 & 0 & -\sin\phi_2 \\
0 & \cos\phi_1 & \sin\phi_1 \cos\phi_2 \\
0 & -\sin\phi_1 & \cos\phi_1 \cos\phi_2
\end{bmatrix}
\begin{bmatrix}
\dot{\theta}_1 \\[2pt] \dot{\theta}_2 \\[2pt] \dot{\theta}_3
\end{bmatrix}.
\label{eq:cinematica_euler}
\end{equation}
A matriz $3\times 3$ em \eqref{eq:cinematica_euler} não é ortogonal,
pois os vetores associados às colunas (no espaço das taxas de Euler) não formam
uma base ortonormal: os versores $\vec{b}_1$, $\vec{a}''_2$ e $\vec{a}'_3$ não
são mutuamente ortogonais. Portanto, essa matriz não representa uma rotação,
mas apenas uma transformação linear entre as taxas de Euler e as componentes de
$\vec{\omega}$.

A relação inversa,
obtida pela inversão explícita da matriz de \eqref{eq:cinematica_euler},
fornece as equações diferenciais cinemáticas para a sequência
$C_1(\phi_1) \leftarrow C_2(\phi_2) \leftarrow C_3(\phi_3)$:
\begin{equation}
\begin{bmatrix}
\dot{\theta}_1 \\[2pt] \dot{\theta}_2 \\[2pt] \dot{\theta}_3
\end{bmatrix}
=
\frac{1}{\cos\phi_2}
\begin{bmatrix}
\cos\phi_2 & \sin\phi_1 \sin\phi_2 & \cos\phi_1 \sin\phi_2 \\
0 & \cos\phi_1 \cos\phi_2 & -\sin\phi_1 \cos\phi_2 \\
0 & \sin\phi_1 & \cos\phi_1
\end{bmatrix}
\begin{bmatrix}
\omega_1 \\[2pt] \omega_2 \\[2pt] \omega_3
\end{bmatrix}.
\label{eq:euler_inversa}
\end{equation}
Conforme \cite{wie2008space}, que denomina \eqref{eq:euler_inversa} de equação cinemática diferencial para esse conjunto de ângulos de Euler.

A partir dos três ângulos $(\phi_1,\phi_2,\phi_3)$, a matriz de cossenos diretores correspondente é dada pela composição das três rotações elementares:

\begin{equation}
\mathbf{C}^{B/A}
=
\mathbf{C}_1(\phi_1)\,\mathbf{C}_2(\phi_2)\,\mathbf{C}_3(\phi_3),
\end{equation}

que é equivalente a uma sequência de rotações extrínsecas em torno dos eixos fixos $x$–$y$–$z$, na ordem $X$–$Y$–$Z$. Essa forma facilita a comparação entre a descrição matricial (DCM) e a parametrização por ângulos de Euler.

A matriz em \eqref{eq:cinematica_euler} torna-se singular quando
$\phi_2 = \pm 90^\circ$, pois $\cos\phi_2 = 0$ e a inversa
\eqref{eq:euler_inversa} deixa de existir. Nessa configuração, dois eixos de
rotação tornam-se colineares, o que reduz efetivamente o número de graus de
liberdade — fenômeno conhecido como singularidade cinemática ou
\emph{gimbal lock}. A escolha de outra sequência de Euler apenas desloca a
singularidade para outra combinação de ângulos; a presença de singularidades é
intrínseca a qualquer parametrização por ângulos de Euler. Essa limitação motiva
o emprego de representações alternativas de atitude, como quaternions, que
evitam essas singularidades no espaço de parâmetros.
